% !TeX root = ./main.tex

% --------------------------------------------------
% 資訊設定（Information Configs）
% --------------------------------------------------

\ntusetup{
  university*   = {National Taiwan University},
  university    = {國立臺灣大學},
  college       = {工學院},
  college*      = {College of Engineering},
  institute     = {工程科學及海洋工程學系},
  institute*    = {Department of Engineering Science and Ocean Engineering},
  title         = {基於標註擴張連結之表皮內神經纖維免疫螢光影像重建演算法},
  title*        = {An Annotation Expansion Linker for Intraepidermal Nerve Fiber Reconstruction in Immunofluorescence Images},
  author        = {溫紹傑},
  author*       = {Shao-Jie Wen},
  ID            = {R13525122},
  advisor       = {張恆華},
  advisor*      = {Herng-Hua Chang},
  date          = {2026-07-30},
  oral-date     = {2026-07-30},
  DOI           = {10.6342/NTU202603123},
  keywords      = {表皮內神經纖維密度、小纖維神經病變、神經纖維重建、多源最短路徑、最小生成森林},
  keywords*     = {Intraepidermal Nerve Fiber Density, Small Fiber Neuropathy, Neural Fiber Reconstruction, Multi-source Shortest Path, Minimum Spanning Forest},
}

% --------------------------------------------------
% 加載套件（Include Packages）
% --------------------------------------------------

\usepackage[numbers,sort&compress]{natbib}      % 參考文獻（abbrv 為數字樣式，須以 numbers 選項配合）
\usepackage{amsmath, amsthm, amssymb}   % 數學環境
\usepackage[normalem]{ulem}             % 下劃線、雙下劃線與波浪紋效果（normalem 保留 \emph 為斜體；xeCJK 不相容 CJKulem）
\usepackage{booktabs}                   % 改善表格設置
\usepackage{multirow}                   % 合併儲存格
\usepackage{diagbox}                    % 插入表格反斜線
\usepackage{array}                      % 調整表格高度
\usepackage{longtable}                  % 支援跨頁長表格
\usepackage{paralist}                   % 列表環境


\usepackage{lipsum}                     % 英文亂字
\usepackage{zhlipsum}                   % 中文亂字
\usepackage{float}                      % 浮動體控制 (H 選項)
\usepackage{algorithm}                  % 演算法環境
\floatname{algorithm}{演算法}           % 將 Algorithm 標題中文化
\usepackage{graphicx}                   % 插圖

% --------------------------------------------------
% 套件設定（Packages Settings）
% --------------------------------------------------


% --------------------------------------------------
% 縮寫列表環境（List of Abbreviations）
% 鏡像 class 之 denotation 環境，僅替換標題；專放英文縮寫，與符號列表分離
% --------------------------------------------------
\makeatletter
\newcommand{\ntu@abbreviation@zh}{縮寫列表}
\newcommand{\ntu@abbreviation@en}{List of Abbreviations}
\ifthenelse{\equal{\ntu@language}{english}}
  {\newcommand{\ntu@abbreviation}{\ntu@abbreviation@en}}
  {\newcommand{\ntu@abbreviation}{\ntu@abbreviation@zh}}

\newenvironment{abbreviations}[1][2.5cm]{%
  \chapter*{\centering \ntu@abbreviation}%
  \addcontentsline{toc}{chapter}{\ntu@abbreviation}%
  \noindent
  \begin{list}{}{%
    \renewcommand\makelabel[1]{##1\hfill}%
    \setlength{\labelwidth}{#1}%                  % 縮寫欄寬度
    \setlength{\labelsep}{0.5cm}%                 % 標籤與列表文字距離
    \setlength{\itemindent}{0cm}%                 % 標籤縮進
    \setlength{\leftmargin}{\labelwidth+\labelsep}% % 標籤左邊界
    \setlength{\rightmargin}{0cm}%                % 標籤右邊界
    \setlength{\parsep}{0cm}%                     % 段落間距
    \setlength{\itemsep}{18pt}%                   % 標籤間距
    \setlength{\listparindent}{0cm}%             % 段落縮進
    \setlength{\topsep}{0pt}%                     % 標籤與上文距離
  }%
}{\end{list}}
\makeatother
